% Aula 10 --- o voto de k vizinhos: a variância cai com k, e o 1-NN erra no máximo o dobro do classificador de Bayes.
Otávio é agrônomo e usa o KNN para classificar os talhões de uma fazenda em atacados por uma
praga ($Y=1$) ou não ($Y=0$), a partir de medidas de solo e de clima. Num talhão novo, com
covariáveis $\x$, a probabilidade de praga é $\eta(\x)=\Prob(Y=1\mid\X=\x)=0{,}6$. O KNN
estima essa probabilidade pela proporção de praga entre os $k$ talhões de treino mais
próximos de $\x$, cujos índices formam o conjunto $\mathcal{N}_{\x}^{(k)}$,
\[
  \widehat P=\frac1k\sum_{i\in\mathcal{N}_{\x}^{(k)}}Y_i ,
\]
e declara praga quando $\widehat P>1/2$. Tome $k$ ímpar, para que não haja empate. Para
isolar o papel de $k$, suponha que os rótulos dos vizinhos sejam independentes entre si e
do rótulo $Y$ do talhão novo, e que cada um seja igual a $1$ com probabilidade $0{,}6$, como
se todos os vizinhos estivessem no próprio $\x$.
\begin{itens}
  \item \pts{1,1} Mostre que $\E[\widehat P]=0{,}6$ e $\Var(\widehat P)=0{,}24/k$.
  \item \pts{1,8} Seja $q_k=\Prob(\widehat P>1/2)$. Mostre que a probabilidade de o KNN errar
        o talhão novo é $0{,}6-0{,}2\,q_k$. Calcule $q_5$ e, usando $q_1=0{,}600$,
        $q_{15}\approx0{,}787$ e $q_{51}\approx0{,}926$, a probabilidade de erro para
        $k=1,5,15$ e $51$. Use a desigualdade de Chebyshev para mostrar que $q_k\to1$ e que a
        probabilidade de erro tende a $0{,}4$. Que classificador erra com probabilidade
        $0{,}4$ nesse talhão?
  \item \pts{1,1} Mostre que, com $k=1$ e uma probabilidade $\eta$ qualquer no lugar de
        $0{,}6$, a probabilidade de erro é $2\eta(1-\eta)$, e que
        $\min(\eta,1-\eta)\le2\eta(1-\eta)\le2\min(\eta,1-\eta)$. Interprete. O que acontece
        no outro extremo, com $k=n$?
\end{itens}

\begin{solucao}
\textbf{(a)} O conjunto $\mathcal{N}_{\x}^{(k)}$ tem $k$ elementos, e cada $Y_i$, com
$i\in\mathcal{N}_{\x}^{(k)}$, é uma Bernoulli$(0{,}6)$, com média $0{,}6$ e variância
$0{,}6\cdot0{,}4=0{,}24$. A esperança sai da linearidade, e a variância, da independência dos
rótulos, que permite somar as variâncias:
\[
  \E[\widehat P]=\frac1k\sum_{i\in\mathcal{N}_{\x}^{(k)}}\E[Y_i]=\frac{k\cdot0{,}6}{k}=0{,}6,
  \qquad
  \Var(\widehat P)=\frac1{k^2}\sum_{i\in\mathcal{N}_{\x}^{(k)}}\Var(Y_i)
  =\frac{k\cdot0{,}24}{k^2}=\frac{0{,}24}{k} .
\]
Nesta idealização, a proporção de vizinhos é um estimador não viesado de $\eta(\x)$, e a sua
variância cai como $1/k$.

\smallskip
\textbf{(b)} O KNN erra se declara praga num talhão sem praga ou deixa de declarar num talhão
com praga. Como a decisão depende só dos rótulos dos vizinhos, que são independentes de $Y$,
\[
  \Prob(\ghat(\x)\ne Y)=q_k\,\Prob(Y=0)+(1-q_k)\,\Prob(Y=1)=0{,}4\,q_k+0{,}6\,(1-q_k)=0{,}6-0{,}2\,q_k .
\]
Com $k=5$, o KNN declara praga quando pelo menos $3$ dos $5$ vizinhos têm praga, e o número
deles é uma binomial$(5;\,0{,}6)$:
\[
  q_5=\binom53(0{,}6)^3(0{,}4)^2+\binom54(0{,}6)^4(0{,}4)+(0{,}6)^5
     =0{,}3456+0{,}2592+0{,}0778\approx0{,}683 .
\]
As probabilidades de erro são $0{,}480$, $0{,}463$, $0{,}443$ e $0{,}415$ para $k=1$, $5$, $15$ e
$51$. Como $k$ é ímpar, $\widehat P\ne1/2$, e o KNN deixa de declarar praga exatamente
quando $\widehat P<1/2$. Esse evento está contido em
$\{\abs{\widehat P-0{,}6}\ge0{,}1\}$, e, pela desigualdade de Chebyshev,
\[
  1-q_k=\Prob\big(\widehat P<1/2\big)\le\Prob\big(\abs{\widehat P-0{,}6}\ge0{,}1\big)
  \le\frac{\Var(\widehat P)}{0{,}1^2}=\frac{24}{k}\ \longrightarrow\ 0 .
\]
Logo $q_k\to1$ e a probabilidade de erro tende a $0{,}6-0{,}2=0{,}4$. É a probabilidade de
erro do classificador de Bayes nesse talhão: como $\eta(\x)>1/2$, ele declara praga, e erra
com probabilidade $1-\eta(\x)=0{,}4$, que nenhum classificador pode melhorar.

\smallskip
\textbf{(c)} Com $k=1$, o KNN repete o rótulo do vizinho mais próximo, que é $1$ com
probabilidade $\eta$ e independente de $Y$. Ele erra se os dois rótulos diferem:
\[
  \Prob(\ghat(\x)\ne Y)=\eta(1-\eta)+(1-\eta)\eta=2\eta(1-\eta) .
\]
Sejam $m=\min(\eta,1-\eta)$ e $1-m=\max(\eta,1-\eta)$, com $1/2\le1-m\le1$. Então
$2\eta(1-\eta)=2m(1-m)$, e multiplicar as desigualdades $1/2\le1-m\le1$ por $2m\ge0$ dá
$m\le2m(1-m)\le2m$. Como $m$ é a probabilidade de erro do classificador de Bayes em $\x$, o
$1$-NN, nesta idealização, erra pelo menos tanto quanto ele e no máximo o dobro; com
$\eta=0{,}6$, são $0{,}48$ contra $0{,}40$. No outro extremo, com $k=n$, todos os pontos de treino
votam, $\widehat P$ é a proporção de talhões com praga no treino inteiro, e o KNN atribui a
mesma classe, a majoritária, a qualquer talhão: é o classificador constante, com o maior viés
possível.
\end{solucao}
